\section{Geometric capture and saddle preparation}
\label{app:geometric-saddle-details}

\subsection{The isolating block and whole-graphic capture}
\label{app:graphic-capture-details}

We prove \cref{prop:isolation-capture-new} in two steps. An outward
transverse boundary around the finite equilibrium gives the isolating
block. Near the graphic, flow boxes force crossings of local section
segments, and intersection numbers in a tubular annulus show that
each section is crossed exactly once.

At the finite equilibrium, the linear part of the organizing field \(\mathcal X_0\) is
\[
 \mathsf L_0=\begin{pmatrix}1&-1\\ C_0&0\end{pmatrix}.
\]
The symmetric matrix
\[
 \mathsf Q_0=\begin{pmatrix}
 (C_0+1)/2&-1/2\\[1mm]
 -1/2&(C_0+2)/(2C_0)
 \end{pmatrix}
\]
is positive definite and satisfies
$\mathsf L_0^T\mathsf Q_0+\mathsf Q_0\mathsf L_0=\mathsf I$.  Hence, for
$\mathsf Q(u)=u^T\mathsf Q_0u$,
\[
 \dot{\mathsf Q}(u)=|u|^2+O(|u|^3),
\]
where the dot denotes the derivative along $\mathcal X_0$.
Choose a sufficiently small disk $B_{\rm eq}=\{\mathsf Q\le\rho_{\rm eq}^2\}$,
disjoint from $\GammaI$, for which the vector field points strictly out of
$B_{\rm eq}$.  Then $B_{\mathrm{iso}}=\mathbb S^2\setminus\operatorname{int}B_{\rm eq}$
has a boundary along which the field points strictly inward. This
property persists in a full parameter neighborhood and proves (i).

The periodic orbits of \eqref{eq:quadratic-family} lie in the open upper
hemisphere of $\mathbb S^2$. The graphic
$\GammaI$ is a smooth simple closed curve in $\mathbb S^2$. Away from the nilpotent
point it is the parabola \eqref{eq:organizing-parabola}. In the chart
\eqref{eq:projective-coordinates-new}, whose coordinate $z$ takes both signs
on the sphere, the part of the parabola with $y>0$ is the graph
\[
 z=\frac{\sqrt{1+C_0v^2}-1}{C_0}=\frac{v^2}{2}+O(v^4),\qquad v\ne0,
\]
and its closure passes smoothly through the nilpotent point $(v,z)=(0,0)$.
Hence $\GammaI$ has closed tubular neighborhoods in
$\operatorname{int}B_{\mathrm{iso}}$; each of them is a compact annulus,
since $\mathbb S^2$ is orientable.

Take \(\Sigma_0\) on the affine orbit and the two projective faces
\(\Sigma_{\rm in},\Sigma_{\rm out}\) on its regular branches.
Choose a closed tubular annulus
\(B_{\rm cap}\Subset\operatorname{int}B_{\rm iso}\) so thin that
\[
 S_j=\Sigma_j\cap B_{\rm cap},\qquad
 j\in\{0,\mathrm{in},\mathrm{out}\},
\]
are disjoint proper spanning arcs in the cyclic order of \(\GammaI\).
Choose a smaller tube
\(\overline{\mathcal N'}\Subset B_{\rm cap}\) for which
\(T_j=S_j\cap\overline{\mathcal N'}\) are compact transverse cores.
Around the marked points \(\xi_j=\GammaI\cap S_j\), choose nested,
pairwise disjoint flow boxes \(F'_j\Subset F_j\Subset B_{\rm cap}\),
with \(T_j\subset F_j\), such that every trajectory through \(F'_j\)
meets \(T_j\) before leaving \(F_j\), in positive or negative time.
All crossings of each \(T_j\) have the organizing orientation.
These properties persist for nearby parameters by transversality
and smooth dependence of the flow boxes.

For sufficiently small \(\delta>0\),
\[
 N_\delta(\GammaI)\subset\mathcal N',\qquad
 N_\delta(\xi_j)\subset F'_j,
\]
where \(N_\delta\) denotes the metric neighborhood. If
\(\distH(\overline\gamma,\GammaI)<\delta\), one Hausdorff inclusion
puts \(\gamma\) in \(\mathcal N'\), and the other supplies a point in
each \(F'_j\). Hence \(\gamma\) crosses every \(T_j\) at least once,
and all its intersections with \(S_j\) lie in that core and have the
same sign.

Let \(\pi:B_{\rm cap}\to\GammaI\) be the tubular projection.
The algebraic intersection with a spanning arc is
\(I(\gamma,S_j)=\deg(\pi|_\gamma)\), up to orientation, and is nonzero.
Thus \(\gamma\), an embedded circle when taken with minimal period,
is essential in the annulus. Such a circle represents a primitive
generator of \(H_1(B_{\rm cap};\mathbb Z)\); its intersection number
with each spanning arc has absolute value one. Since all crossings
have the same sign, there is exactly one per arc. Disjoint spanning
arcs are met in their cyclic order, and the flow fixes that order to
be the direction of \(\GammaI\). This proves (ii).

\subsection{Uniform preparation of the interior saddle}
\label{app:positive-saddle-preparation}

For \eqref{eq:direct-positive-scaled-field}, we first straighten the
stable and unstable graphs with uniform finite-order bounds, then
normalize the equations on the two axes. The parameter-dependent
stable manifold theorem \cite[Theorem~9.6]{Teschl2012} supplies the
graphs through the following amplitude comparison.

Fix a finite order \(k\), with extra derivatives for the coordinate
changes below. In the scaled parameters
\(\widehat p=(A,C,\kappa,q_2,\ell)\), where
\(q_0=\ell^2\) and \(q_1=\kappa\ell\), the field extends smoothly
to a neighborhood of a compact parameter box including \(\ell=0\).
Factoring \(B(\ell X,\ell^2S,p)-B(0,0,p)\) by \(\ell\) writes it as
\[
 \mathcal X_{\widehat p}=\mathcal X^0_{B_0(p)}
   +\kappa\mathcal X_1+q_2\mathcal X_2+\ell\mathcal X_{3,\widehat p},
 \qquad \mathcal X_1=(X,0),\quad \mathcal X_2=(X^2,0),
\]
where \(\mathcal X^0_{B_0}\) is \eqref{eq:direct-positive-comparison} and
\(\mathcal X_{3,\widehat p}\) is smooth with uniformly bounded state derivatives on
the fixed state box. Replace the displayed coefficients
\(\kappa,q_2,\ell\) by independent small amplitudes
\(\alpha=(\alpha_1,\alpha_2,\alpha_3)\), keeping \(\widehat p\) fixed.
At \(\alpha=0\), the comparison
saddle has eigenvalues \(1,-1\); reducing the amplitude box gives
a common positive spectral margin. The cited theorem on a finite
cover of the compact \((\widehat p,\alpha)\) set gives stable and unstable
graphs on a common smaller state box, with bounded state derivatives
and one additional amplitude derivative. The mean value theorem from
\(\alpha=0\) to \(\alpha=(\kappa,q_2,\ell)\) gives graph differences
\(O_{C^k}(\kappa+|q_2|+\ell)\). Smooth eigen-coordinates centered
at each saddle have the same bounds; straightening the graphs gives
\[
 x'=xg_{\rm u}(x,y),\qquad y'=-yg_{\rm s}(x,y),\qquad g_{\rm u},g_{\rm s}>c>0.
\]
An orbital time change makes \(x'=x,\ y'=-yg(x,y)\),
where \(g(0,0)=\nu_p\).

Normalize the \(x=0\) equation by the regular coordinate
\[
 \psi(y)=y\exp\int_0^y
           \left(\frac{\nu_p}{g(0,t)}-1\right)\frac{dt}{t}.
\]
The new equation is \(y'=-y[\nu_p+xc_1(x,y)]\).
Multiplication of \(y\) by
\(\exp\int_0^xc_1(t,0)\,dt\) then gives exactly
the prepared equation \eqref{eq:direct-positive-saddle-preparation}.
